%% Assignment 1 - Reactive Following with Five Robots (Webots + ROS 2)
%% IEEE conference format, non-blind as required by the assignment brief.
%% Build:  pdflatex reactive_following.tex   (twice, for references)
\documentclass[conference,a4paper]{IEEEtran}
\usepackage{cite}
\usepackage{amsmath}
\usepackage{graphicx}
\usepackage{booktabs}
\usepackage{url}
\usepackage{siunitx}
\sisetup{detect-all}

\begin{document}

\title{Reactive Following with Five Robots:\\
A Five-Rule Controller in Webots and ROS~2}

\author{
\IEEEauthorblockN{Lara Dehandschutter}
\IEEEauthorblockA{FEUP / Universidade do Porto\\
lara.dehandschutter@ugent.be}
\and
\IEEEauthorblockN{Lennert De Rouck}
\IEEEauthorblockA{FEUP / Universidade do Porto\\
lennert.derouck@ugent.be}
\and
\IEEEauthorblockN{Pedro Cancela}
\IEEEauthorblockA{FEUP / Universidade do Porto\\
up202400230@up.pt}
}

\maketitle

\begin{abstract}
We present a purely reactive controller that makes five differential-drive
robots form a following chain in Webots, driven by five ROS~2 nodes. Each
robot acts on a single current snapshot from a frontal camera and LiDAR
($100^\circ$, \SI{4}{\metre}); there is no target memory, no global pose and no
inter-robot communication. Robot orientation is read directly from colour:
the front half of each chassis is red and the rear half green, so ``facing
me'' and ``facing away'' reduce to which colour is nearer. The controller is
five prioritised rules. In the reported run the five robots formed a stable
train --- one leader on its base curve and four followers --- after
\SI{186}{\second}, and held it for the remaining \SI{105}{\second} with no
avoidance event of any kind. We report following duration, angular error and
target losses, and devote the larger part of this paper to the quirks and
failure modes that remain, several of which are consequences of the sensing
envelope rather than of the control rules.
\end{abstract}

\section{Introduction}
The task is reactive following: five robots, each deciding only from its own
current sensor readings, should aggregate into one or more following chains.
The constraints are those of classical reactive control in the sense of
Brooks' subsumption architecture~\cite{brooks1986}: behaviour is a priority
ordering over situation--action rules, with no world model and no planning.
The assignment additionally forbids GPS, global localisation, target memory
and privileged state, while permitting internal controller state of the
PI/PID kind.

Our contribution is deliberately small. The interesting content of this report
is not the controller, which is roughly 120 lines, but the catalogue of ways
in which a correct-looking reactive rule set misbehaves once five copies of it
interact in a shared environment. Section~\ref{sec:limitations} is therefore
the longest section, by intent.

\section{Simulation Model and Sensing}

\subsection{Robot}
Each robot is a differential-drive cylinder of radius \SI{0.115}{\metre} with
two powered wheels and an invisible contact bumper of radius
\SI{0.146}{\metre}. Two robots are in contact when their centres are
\SI{0.292}{\metre} apart; this number is the reference against which all
proximity results below are read.

\subsection{Colour-coded orientation}
An earlier iteration used two small recognition plaques, one red at the front
and one green at the rear. These were replaced by two coloured half-shells
covering the front and rear halves of the chassis. The reason is angular
coverage: a flat plaque shrinks to nothing in projection as a robot turns
away, and we measured that the leader became unclassifiable once the angle
between its heading and the observer's line of sight exceeded about
$74^\circ$ --- precisely when a leader was turning, which is when a follower
most needs to keep it.

Each half is a separate Webots \texttt{Solid} with its own
\texttt{recognitionColors}. One detail is load-bearing: Webots reports the
\emph{origin} of a recognised Solid, so the two halves are offset to
$x=\pm\SI{0.060}{\metre}$. Sharing the chassis centre would make both halves
report the same position and front/back would be indistinguishable.

\subsection{Sensors}
Sensing is frontal only, as required: a recognition camera and a LiDAR, both
$100^\circ$ and \SI{4}{\metre}. The LiDAR mounting deserves comment because it
caused most of our early failures and its fix is non-obvious.

Mounted centrally on the roof at $z=\SI{0.155}{\metre}$ --- the natural place,
and necessary to clear its own chassis --- the LiDAR sits \emph{above} the
\SI{0.125}{\metre} body top of every other robot. With a narrow vertical field
of view, the beam only descends to body height at
\begin{equation}
d \;=\; \frac{0.155-0.125}{\tan(\theta_v/2)} \;\approx\; \SI{0.75}{\metre},
\end{equation}
so a peer closer than \SI{0.75}{\metre} is invisible. Every collision
threshold we might wish to set lies inside that blind radius. Worse, the
blindness is unavoidable for a \emph{central} sensor: to avoid seeing its own
rim at radius $r_b$ the beam drop at $r_b$ must stay under the clearance $c$,
while seeing a peer at distance $d$ requires the drop at $d-r_b$ to reach $c$.
Both $\theta_v$ and $c$ cancel, leaving $d > 2r_b$ --- twice the body radius,
which is exactly the distance at which two robots already touch.

The resolution exploits the fact that the LiDAR is now frontal. Mounted at the
\emph{front edge} ($x=\SI{0.130}{\metre}$, outside the body radius), every
forward ray's closest approach to the chassis lies behind the sensor, so it
can be lowered to $z=\SI{0.115}{\metre}$: inside the \SIrange{0.045}{0.125}{\metre}
band of every peer. Robots are then visible at all ranges including contact.
A $360^\circ$ LiDAR cannot do this, because its rearward rays would face its
own chassis.

\subsection{Software architecture}
Webots runs one controller per robot, which packs a camera/LiDAR snapshot into
a UDP datagram. A ROS~2 bridge republishes each as
\texttt{/R\textit{i}/local\_sensors}; five independent \texttt{ReactiveFollower}
nodes subscribe and publish \texttt{/R\textit{i}/cmd\_vel}. No node subscribes
to another robot's topics, which is what makes the no-communication constraint
auditable rather than merely asserted.

\section{Reactive Controller}
Perception reduces each frame to the nearest red and the nearest green, as
(distance, bearing) pairs. The controller is five rules in priority order
(Table~\ref{tab:rules}).

\begin{table}[t]
\caption{Rules, in evaluation order. First match wins.}
\label{tab:rules}
\centering
\small
\begin{tabular}{@{}clll@{}}
\toprule
\# & State & Condition & Action \\
\midrule
1 & \texttt{ESCAPE\_TURN} & red nearer, $d_r<0.45$ & $v=0$, turn right \\
2 & \texttt{AVOID\_RED}   & red nearer, $d_r<0.90$ & turn right \\
3 & \texttt{FOLLOW}       & green visible          & steer to it, speed ramp \\
4 & \texttt{APPROACH}     & red only, $d_r\ge0.90$ & steer to it \\
5 & \texttt{AVOID\_WALL}  & no robot, LiDAR $<2.0$ & turn to open side \\
-- & \texttt{BASE\_CURVE} & otherwise              & $v_0$, $\omega_{\text{base}}$ \\
\bottomrule
\end{tabular}
\end{table}

Two design points proved necessary rather than cosmetic.

\paragraph{Red is a comparison, not a detection}
Rules 1--2 require red to be \emph{nearer than green}, not merely visible.
Every robot being followed presents a red half somewhere in view, so keying on
visibility made followers abandon correct leaders.

\paragraph{Turn rate is capped by radius, not by a fixed rate}
Steering straight at a leader is pure pursuit, which cuts inside the leader's
path. In a five-robot ring the leader sits about \SI{0.63}{\radian} off-axis,
giving $\omega = 1.5\times0.63 = \SI{0.94}{\radian\per\second}$ and an arc
radius of $v/\omega \approx \SI{0.15}{\metre}$ against a \SI{0.95}{\metre}
ring: the formation spirals inward until it piles up. We cap by
$\omega_{\max} = v_0/R_{\min}$ rather than by a shared $\omega$, because the
robots run at different speeds (\SIrange{0.12}{0.16}{\metre\per\second}) and a
shared cap would give each a different radius and distort the ring.

\section{Experimental Method}
The five robots start evenly spaced on a \SI{0.95}{\metre} circle, each
heading along the tangent. We record per robot: following duration, mean
absolute bearing error while following, target losses (transitions out of a
following state) and bumper contacts. A seeded script
(\texttt{tools/randomize\_starts.py}) regenerates scattered starting poses so
that runs can be repeated with different positions and headings while
remaining reproducible.

\section{Results}
Table~\ref{tab:results} reports a \SI{291}{\second} run. The controller issued
only 48 mode changes across all five robots, i.e.\ \SI{9.9}{\per\minute}.

\begin{table}[t]
\caption{Per-robot results, \SI{291.2}{\second} run.}
\label{tab:results}
\centering
\small
\begin{tabular}{@{}lrrrr@{}}
\toprule
 & Following & \% & Losses & Mean $|\beta|$ (rad) \\
\midrule
R1 & \SI{290.9}{\second} & 99.9 & 2 & 0.629 \\
R2 & \SI{291.1}{\second} & 99.9 & 1 & 0.702 \\
R3 & \SI{44.7}{\second}  & 15.3 & 2 & 0.420 \\
R4 & \SI{282.8}{\second} & 97.1 & 3 & 0.729 \\
R5 & \SI{137.2}{\second} & 47.1 & 6 & 0.451 \\
\bottomrule
\end{tabular}
\end{table}

Time-weighted occupancy was \texttt{FOLLOW} 71.9\%, \texttt{BASE\_CURVE}
27.5\%, and all avoidance states together below 1\%
(\texttt{AVOID\_RED} 0.3\%, \texttt{ESCAPE\_TURN} 0.2\%,
\texttt{AVOID\_WALL} 0.1\%).

The run passed through five phases: \SI{82}{\second} of scattered search;
\SI{34}{\second} of first acquisitions; an \SI{11}{\second} congestion
incident; \SI{59}{\second} of re-forming; and then \SI{105}{\second} of stable
train with no avoidance event at all. The end state is the intended one --- R3
leading on its base curve, R1, R2, R4 and R5 following --- which is why R3's
low following percentage is a success rather than a failure: the leader of a
train is, by definition, not following anything.

Mean absolute bearing error while following was \SI{0.549}{\radian}
($31.5^\circ$). This should not be read as tracking error. In a five-robot
ring the leader is geometrically about $36^\circ$ off-axis, so the measured
value is close to the best achievable for this formation.

\section{Quirks, Shortcomings and Limitations}
\label{sec:limitations}

\subsection{Edge-of-cone bistability}
The one congestion incident is instructive. R3 and R5 each detected an
oncoming robot, turned right and cleared in about \SI{1.4}{\second}: the rule
working as designed. R4 instead oscillated between \texttt{AVOID\_WALL} and
\texttt{ESCAPE\_TURN} three times in \SI{4}{\second}. The cause is visible in
the bearing at resolution, $+49.2^\circ$: its target sat on the $\pm50^\circ$
cone edge and flickered in and out of detection. The two states require
opposite readings --- ``red very close'' and ``no robot at all'' --- so the
controller alternated between them. This is a sensor-envelope artefact, not a
logic error; no rule is wrong. A hysteresis band would suppress it at the cost
of a small amount of remembered state.

\subsection{The emergency response is rate-limited}
An acceleration limiter smooths all commands. It also smooths the emergency:
\texttt{ESCAPE\_TURN} commands $v=0,\ \omega=-0.80$, but the log shows
$v=0.12,\ \omega=+0.37$ at onset, i.e.\ the robot is still rolling forward and
still turning \emph{left} for roughly \SI{0.2}{\second} while the limiter
unwinds the previous command. The highest-priority rule is therefore not
actually the fastest-acting one.

\subsection{Collisions were not captured}
The metrics summary prints only on a clean shutdown. The analysed run was
killed rather than interrupted, so bumper counts were lost. From mode
thresholds we can infer that \texttt{ESCAPE\_TURN} triggers at
\SI{0.64}{\metre} centre-to-centre against a \SI{0.292}{\metre} contact
distance, so the incident was clear at the trigger --- but the logs of that
run cannot say how much closer it came. Distance logging and a
closest-approach metric were added afterwards and are not reflected in the
numbers above. We report this rather than quietly omitting the metric.

\subsection{No rear or lateral sensing}
A robot alongside or behind is invisible: the camera and LiDAR both span only
$100^\circ$ forward. Sideways collisions are therefore undetectable in
principle, not merely unhandled. Our mitigations are preventive --- lane
geometry and turn-rate caps that stop robots ending up abreast --- and nothing
reacts once they are.

\subsection{Open-loop geometric coupling}
Three parameters are not independent. For five robots on a ring, spacing and
radius satisfy $s = 2R\sin(\pi/5) = 1.176R$, and the camera reads
\SI{0.19}{\metre} less than centre spacing. An earlier standoff of
\SI{0.45}{\metre} therefore demanded a \SI{0.54}{\metre} ring while the turn
cap refused anything tighter than \SI{0.90}{\metre}; the robots bunched and
crawled at the speed floor. The present values agree, but they were made to
agree by hand and nothing in the code enforces the relationship.

\subsection{Mode changes are the only log events}
Modes are logged on change, so a quantity that drifts \emph{within} a mode is
invisible. We initially misdiagnosed a collision because \texttt{APPROACH} was
recorded at $\omega=0.27$ while its true peak, reached between log lines, was
at least $0.73$. Any conclusion drawn from this log format understates
within-state extremes.

\subsection{Measurement artefacts}
One earlier analysis attributed 70\% of robot-time to sensor dropouts and
concluded the simulation was running below real time. The log had in fact been
recorded while stepping Webots tick by tick, so the simulation was halted
between clicks. When run freely, \texttt{SENSOR\_FAILSAFE} occupies 0.0\% of
the run. We record the mistake because the artefact is easy to reproduce and
the wrong conclusion was superficially well-supported.

\subsection{Repository state}
The test suite does not match the delivered controller. All ten test modules
exercise the previous \texttt{safe\_logic} cascade and the marker-based
perception module, neither of which is wired into the current system; they
would not run against it. The old controller is retained in the repository for
comparison but is not executed.

\subsection{Single-run evidence}
The results are one run. The seeded randomiser exists so that repeats with
different initial conditions can be performed, and the comparison across
turning strengths that the brief suggests (no following / weak / strong) is
supported by a command-line flag, but neither has been run at the time of
writing. No statistical claim is made.

\section{Conclusion}
A five-rule reactive controller, acting on one frontal snapshot at a time,
assembled five robots into a stable leader-plus-four-followers train in
\SI{186}{\second} and held it without further intervention. The rules were
cheap; the difficulty was almost entirely in the interaction between the
sensing envelope and the geometry of a multi-robot formation. The dominant
failure mode we did not eliminate --- bistable detection at the edge of the
field of view --- is a property of where the sensor stops, not of what the
controller decides.

\begin{thebibliography}{9}
\bibitem{brooks1986}
R.~A. Brooks, ``A robust layered control system for a mobile robot,''
\emph{IEEE J. Robot. Autom.}, vol.~2, no.~1, pp.~14--23, 1986.

\bibitem{arkin1998}
R.~C. Arkin, \emph{Behavior-Based Robotics}. Cambridge, MA, USA: MIT Press,
1998.

\bibitem{webots}
O.~Michel, ``Webots: Professional mobile robot simulation,''
\emph{Int. J. Adv. Robot. Syst.}, vol.~1, no.~1, pp.~39--42, 2004.

\bibitem{ros2}
S.~Macenski, T.~Foote, B.~Gerkey, C.~Lalancette, and W.~Woodall,
``Robot Operating System 2: Design, architecture, and uses in the wild,''
\emph{Science Robotics}, vol.~7, no.~66, 2022.

\bibitem{coulter1992}
R.~C. Coulter, ``Implementation of the pure pursuit path tracking algorithm,''
Robotics Institute, Carnegie Mellon Univ., Tech. Rep. CMU-RI-TR-92-01, 1992.
\end{thebibliography}

\end{document}
